% GENERATED FILE. Edit canonical lessons, then run npm run generate:books.
% canonical-source-hash: fnv1a64:8df4cc8291c7f8cb
% canonical-lessons: SA-C329-puja, SA-C329-ghanta, SA-C329-dhupah, SA-C329-aratrikam, SA-C329-prarthayate

\chapter{At the Temple}
\label{ch:sa-at-the-temple}
\hlchaptermodality{329}

\begin{chapteropening}
\textbf{By the end of this chapter:} \emph{I can say what people do at the temple.}

Complete the last lesson of chapter 329: I can say what people do at the temple.
\end{chapteropening}

\section[\texorpdfstring{\sk{पूजा} (pūjā)}{पूजा (pūjā)}]{\sk{पूजा} (pūjā) — worship}

\label{lesson:SA-C329-puja}



\begin{quote}
Before the new one: say the Sanskrit for offers, then the Sanskrit for a blessed offering.
\end{quote}

\subsection*{You'll want to know: \sk{पूजा}}
\textbf{\sk{पूजा}} — \emph{pūjā} — \textquotedblleft{}worship\textquotedblright{}.

\begin{quote}
\sk{प्रातः} \sk{मन्दिरे} \sk{पूजा} \sk{भवति।}
\end{quote}

\begin{quote}
\emph{prātaḥ mandire pūjā bhavati.}
\end{quote}

\begin{quote}
There is worship in the temple in the morning.
\end{quote}

A feminine noun.

\begin{cousinweb}
\textbf{\sk{पूजा}} (\emph{pūjā}) comes from the root √pūj, \textquotedblleft{}to honour, worship\textquotedblright{}, as in \emph{pūjayati}; MW gives \emph{pūjā} as honour, worship, respect, reverence (MW).
\end{cousinweb}

\begin{grammarlens}[title={What you've built}]
One more word for telling what happened.
\end{grammarlens}

\subsection*{Your turn}
\begin{itemize}
\raggedright
  \item \emph{Say it:} \emph{pūjā}
  \item \emph{Say it:} \emph{pūjā}, once more
  \item \emph{Say it:} \emph{prātaḥ mandire pūjā bhavati.} — There is worship in the temple in the morning.
  \item \emph{Recall:} say the Sanskrit for offers, then the Sanskrit for a blessed offering, then say \emph{pūjā} again
\end{itemize}

\begin{tcolorbox}[breakable,colback=teal!4,colframe=teal!35!black,title={Before you move on}]
What does \sk{पूजा} mean? (\textquotedblleft{}Worship\textquotedblright{}.) Say it once more.
\end{tcolorbox}
\section[\texorpdfstring{\sk{घण्टा} (ghaṇṭā)}{घण्टा (ghaṇṭā)}]{\sk{घण्टा} (ghaṇṭā) — a bell}

\label{lesson:SA-C329-ghanta}



\begin{quote}
Before the new one: say the Sanskrit for a blessed offering, then the Sanskrit for worship.
\end{quote}

\subsection*{You'll want to know: \sk{घण्टा}}
\textbf{\sk{घण्टा}} — \emph{ghaṇṭā} — \textquotedblleft{}a bell\textquotedblright{}.

\begin{quote}
\sk{मन्दिरस्य} \sk{घण्टायाः} \sk{शब्दः} \sk{मधुरः} \sk{अस्ति।}
\end{quote}

\begin{quote}
\emph{mandirasya ghaṇṭāyāḥ śabdaḥ madhuraḥ asti.}
\end{quote}

\begin{quote}
The sound of the temple bell is sweet.
\end{quote}

A feminine noun.

\begin{cousinweb}
\textbf{\sk{घण्टा}} (\emph{ghaṇṭā}) MW gives \emph{ghaṇṭā} as a bell, and a metal plate struck as a gong, without a derivation, so its origin is uncertain (MW).
\end{cousinweb}

\begin{grammarlens}[title={What you've built}]
One more word for telling what happened.
\end{grammarlens}

\subsection*{Your turn}
\begin{itemize}
\raggedright
  \item \emph{Say it:} \emph{ghaṇṭā}
  \item \emph{Say it:} \emph{ghaṇṭā}, once more
  \item \emph{Say it:} \emph{mandirasya ghaṇṭāyāḥ śabdaḥ madhuraḥ asti.} — The sound of the temple bell is sweet.
  \item \emph{Recall:} say the Sanskrit for a blessed offering, then the Sanskrit for worship, then say \emph{ghaṇṭā} again
\end{itemize}

\begin{tcolorbox}[breakable,colback=teal!4,colframe=teal!35!black,title={Before you move on}]
What does \sk{घण्टा} mean? (\textquotedblleft{}A bell\textquotedblright{}.) Say it once more.
\end{tcolorbox}
\section[\texorpdfstring{\sk{धूपः} (dhūpaḥ)}{धूपः (dhūpaḥ)}]{\sk{धूपः} (dhūpaḥ) — incense}

\label{lesson:SA-C329-dhupah}



\begin{quote}
Before the new one: say the Sanskrit for worship, then the Sanskrit for a bell.
\end{quote}

\subsection*{You'll want to know: \sk{धूपः}}
\textbf{\sk{धूपः}} — \emph{dhūpaḥ} — \textquotedblleft{}incense\textquotedblright{}.

\begin{quote}
\sk{मन्दिरं} \sk{धूपेन} \sk{पूर्णम्} \sk{अस्ति।}
\end{quote}

\begin{quote}
\emph{mandiraṃ dhūpena pūrṇam asti.}
\end{quote}

\begin{quote}
The temple is full of incense.
\end{quote}

A masculine noun.

\begin{cousinweb}
\textbf{\sk{धूपः}} (\emph{dhūpaḥ}) comes from the root √dhū, as MW explains it, comparing \emph{puṣpa} from √puṣ; MW gives \emph{dhūpa} as incense, aromatic smoke (MW).
\end{cousinweb}

\begin{grammarlens}[title={What you've built}]
One more word for telling what happened.
\end{grammarlens}

\subsection*{Your turn}
\begin{itemize}
\raggedright
  \item \emph{Say it:} \emph{dhūpaḥ}
  \item \emph{Say it:} \emph{dhūpaḥ}, once more
  \item \emph{Say it:} \emph{mandiraṃ dhūpena pūrṇam asti.} — The temple is full of incense.
  \item \emph{Recall:} say the Sanskrit for worship, then the Sanskrit for a bell, then say \emph{dhūpaḥ} again
\end{itemize}

\begin{tcolorbox}[breakable,colback=teal!4,colframe=teal!35!black,title={Before you move on}]
What does \sk{धूपः} mean? (\textquotedblleft{}Incense\textquotedblright{}.) Say it once more.
\end{tcolorbox}
\section[\texorpdfstring{\sk{आरात्रिकम्} (ārātrikam)}{आरात्रिकम् (ārātrikam)}]{\sk{आरात्रिकम्} (ārātrikam) — the waving of lamps in worship}

\label{lesson:SA-C329-aratrikam}



\begin{quote}
Before the new one: say the Sanskrit for a bell, then the Sanskrit for incense.
\end{quote}

\subsection*{You'll want to know: \sk{आरात्रिकम्}}
\textbf{\sk{आरात्रिकम्}} — \emph{ārātrikam} — \textquotedblleft{}the waving of lamps in worship\textquotedblright{}.

\begin{quote}
\sk{सायं} \sk{मन्दिरे} \sk{आरात्रिकं} \sk{भवति।}
\end{quote}

\begin{quote}
\emph{sāyaṃ mandire ārātrikaṃ bhavati.}
\end{quote}

\begin{quote}
In the evening there is the waving of lamps in the temple.
\end{quote}

A neuter noun.

\begin{cousinweb}
\textbf{\sk{आरात्रिकम्}} (\emph{ārātrikam}) is built on \emph{ā} and \emph{rātri}, \textquotedblleft{}night\textquotedblright{}; MW gives \emph{ārātrika} as the light waved at night before an image, and the ceremony itself (MW).
\end{cousinweb}

\begin{grammarlens}[title={What you've built}]
One more word for telling what happened.
\end{grammarlens}

\subsection*{Your turn}
\begin{itemize}
\raggedright
  \item \emph{Say it:} \emph{ārātrikam}
  \item \emph{Say it:} \emph{ārātrikam}, once more
  \item \emph{Say it:} \emph{sāyaṃ mandire ārātrikaṃ bhavati.} — In the evening there is the waving of lamps in the temple.
  \item \emph{Recall:} say the Sanskrit for a bell, then the Sanskrit for incense, then say \emph{ārātrikam} again
\end{itemize}

\begin{tcolorbox}[breakable,colback=teal!4,colframe=teal!35!black,title={Before you move on}]
What does \sk{आरात्रिकम्} mean? (\textquotedblleft{}The waving of lamps in worship\textquotedblright{}.) Say it once more.
\end{tcolorbox}
\section[\texorpdfstring{\sk{प्रार्थयते} (prārthayate)}{प्रार्थयते (prārthayate)}]{\sk{प्रार्थयते} (prārthayate) — prays for; asks for}

\label{lesson:SA-C329-prarthayate}



\begin{quote}
Before the new one: say the Sanskrit for incense, then the Sanskrit for the waving of lamps.
\end{quote}

\subsection*{You'll want to know: \sk{प्रार्थयते}}
\textbf{\sk{प्रार्थयते}} — \emph{prārthayate} — \textquotedblleft{}prays for; asks for\textquotedblright{}.

\begin{quote}
\sk{वृद्धा} \sk{माता} \sk{शान्तिं} \sk{प्रार्थयते।}
\end{quote}

\begin{quote}
\emph{vṛddhā mātā śāntiṃ prārthayate.}
\end{quote}

\begin{quote}
The old mother prays for peace.
\end{quote}

Past (she …): \sk{सा} \sk{शान्तिं} \sk{प्रार्थितवती}, \emph{sā śāntiṃ prārthitavatī}.

\begin{cousinweb}
\textbf{\sk{प्रार्थयते}} (\emph{prārthayate}) is built from \emph{pra} and the root √arth, \textquotedblleft{}to ask, wish\textquotedblright{}; MW gives \emph{prārthayate} as to wish for, desire, and to ask a person for something (MW).
\end{cousinweb}

\begin{grammarlens}[title={What you've built}]
One more word for telling what happened.
\end{grammarlens}

\subsection*{Your turn}
\begin{itemize}
\raggedright
  \item \emph{Say it:} \emph{prārthayate}
  \item \emph{Say it:} \emph{prārthayate}, once more
  \item \emph{Say it:} \emph{vṛddhā mātā śāntiṃ prārthayate.} — The old mother prays for peace.
  \item \emph{Recall:} say the Sanskrit for incense, then the Sanskrit for the waving of lamps, then say \emph{prārthayate} again
\end{itemize}

\begin{tcolorbox}[breakable,colback=teal!4,colframe=teal!35!black,title={Before you move on}]
What does \sk{प्रार्थयते} mean? (\textquotedblleft{}Prays for; asks for\textquotedblright{}.) Say it once more.
\end{tcolorbox}
